% Generated draft: requires mathematical and source review.
\documentclass[11pt]{article}
\usepackage[margin=25mm]{geometry}
\usepackage{fontspec}
\setmainfont{Latin Modern Roman}
\usepackage{amsmath,amssymb,mathtools,unicode-math}
\setmathfont{Latin Modern Math}
\usepackage{enumitem,needspace,xurl,hyperref}
\setlength{\parindent}{0pt}
\setlength{\parskip}{5pt}
\setlength{\emergencystretch}{3em}
\newcommand{\N}{\mathbb{N}}
\newcommand{\Z}{\mathbb{Z}}
\newcommand{\Q}{\mathbb{Q}}
\newcommand{\R}{\mathbb{R}}
\newcommand{\C}{\mathbb{C}}
\newcommand{\F}{\mathbb{F}}
\newcommand{\A}{\mathbb{A}}
\newcommand{\PP}{\mathbb{P}}
\newcommand{\E}{\mathbb{E}}
\newcommand{\eps}{\varepsilon}
\newcommand{\dd}{\,\mathrm{d}}
\begin{document}
\section*{UnsolvedMath ID: 20001895}
\textbf{Source problem code:} catalogue entry
\subsection{Definitions and mathematical statement}
The extracted text is only a transition: ``Question 5.4.1 above asks whether there are other examples. 5.6. Infinitely many simple closed geodesics.'' The supplied verification says the first sentence ends Section5.5 by referring back to a different question, while the second is a heading. Neither the antecedent examples nor the next numbered question are supplied. This record therefore contains no standalone mathematical assertion to prove, and the heading must not be converted into an invented existence conjecture.

The aid gives the following independently labelled lemma prompted by the heading. Let \(M\) be a connected Riemannian manifold with fundamental group generated by \(r\) elements, let \(D\geq1\) be an integer, and let \(\gamma_1,\ldots,\gamma_N\) be geometrically distinct primitive closed geodesics. Give every covering manifold the pulled-back metric. Assume each \(\gamma_i\) has a simple closed lift in some connected cover of degree at most \(D\). One precise allowed interpretation is an embedded closed geodesic elevation whose projection is the given geodesic image; if one insists that it project with degree one, that stronger hypothesis may be imposed instead. The counting conclusion applies with either convention once fixed consistently.

Put
\[
 C_r(D)=\sum_{d=1}^D(d!)^r.
\]
The reported conclusion is that one connected cover of degree at most \(D\) contains distinct simple lifts of at least
\[
 \left\lceil\frac{N}{C_r(D)}\right\rceil
\]
of the geodesics. The bound uses that degree-\(d\) connected covers are represented by transitive permutation actions of the \(r\)-generated fundamental group, with at most \((d!)^r\) possible homomorphisms into \(S_d\). Distinct projected geodesic images remain distinct upstairs.

Consequently an infinite family with a uniform bound on its simple-lift cover degrees has infinitely many simple lifts in one fixed finite cover. This does not prove that every primitive geodesic lifts simply, supply a uniform \(D\), or answer the missing Question5.4.1 or the next section's unspecified problem. It is an auxiliary conditional result, separate from the context-only source fragment.
\subsection{Short English statement}
This record is a cross-reference plus a section heading, not a standalone problem. The aid's separate lemma says that uniformly bounded simple-lift cover degrees force an infinite geodesic family to concentrate in one finite cover. It supplies no missing existence statement or uniform degree bound.
\subsection{Sources}
\begin{itemize}
\item Local UnsolvedMath registry, permanent ID \(20001895\), code catalogue entry; canonical \texttt{aim-geometry-notes.json}, \texttt{notes[232]}; source tag ``question.''
\item AIM geodesics problem list, transition from Section5.5 to Section5.6, referencing Question5.4.1, \url{https://aimath.org/pastworkshops/geodesicsproblems.pdf}. The supplied verification identifies the fragment boundaries; no PDF inspection is claimed here.
\item No separate authoritative citation is supplied for the bounded-degree cover-counting lemma.
\item Research attempt 1: result ``context\_only,'' problem status at run ``not\_a\_problem,'' registry status ``open,'' rigor label ``proved,'' novelty confidence ``low,'' classification ``OPEN-TRIAGE,'' suggested difficulty L2. These distinguish an auxiliary result from the source fragment. The machine-generated aid is not peer reviewed; independently verify before citation.
\end{itemize}
\end{document}
